% !TEX root = ../print.tex
% Draft \S6 --- Scale covariance: the similarity form
%
% Source: draft/hemigroup-causal-scale-space-kernels.md, lines 161-196.
% Chapter 6 of the blueprint; the shared counter puts every number here on the draft's own.
%
% The chapter that turns "there is some action of the scaling group" into a coordinate. Its
% two rigidity results are proved in the draft as single dense paragraphs -- Lemma 6.2 argues
% four numbered claims in one, with the key monotonicity step inside an em-dash aside -- and
% they are broken out here, statement by statement, with no change of content. That is the
% one place in this chapter where the blueprint should read differently from the draft: this
% is the argument a formalisation has to follow line by line, and the draft's compression is
% the enemy of exactly that. The formalisation note singles out this gauge construction as
% the most unglamorous stretch of the Lean ladder (its M5); the proofs below are written to
% be followed rather than admired.
%
% Equations (6.1) and (6.2) keep the draft's hand-written tags, since its later prose refers
% to them by number.
%
% VOCABULARY (2026-08-10), the companion to the note at the head of 05-cascade.tex. Only one
% conclusion in this chapter names a class of functions -- prop:canonical-gauge's F in BF_0 --
% and it is now also stated in LE (def:levy-exponent), which is what a Lean tag can point at.
%
% The proofs are a second matter and a smaller one. Two steps here reach for the derivative-sign
% picture where they do not have to: lem:action-rigidity(3) argues that a nonzero member of BF_0
% is strictly increasing "because its derivative is completely monotone", and the same fact is
% used again in prop:canonical-gauge. Both are cor:strict-monotonicity, whose Lean route runs
% lem:vanishing on the representation and never differentiates -- see that node. The prose is
% left as the draft has it, since the paper argues in BF_0; the formalisation notes say which
% substitute to use, so that a Lean pass does not rediscover the question mid-proof.

% draft: Lemma 6.1
\begin{lemma}[Laplace form of covariance]\label{lem:covariance-laplace}
\uses{def:cascade-family,lem:convolution-representation,lem:additivity,prop:laplace-uniqueness}
\lean{Hemigroup.CascadeCore.covariance_laplace}\leanok
Under (A1)--(A5), axiom (A8) is equivalent to
\[
  \dil_\sigma\,\mu_{x,y} = \mu_{\Sact_\sigma x,\, \Sact_\sigma y}
  \qquad \Longleftrightarrow \qquad
  g_{\Sact_\sigma x,\, \Sact_\sigma y}(s) = g_{x,y}(\sigma s),
\]
and in particular
\[
  G(\Sact_\sigma x,\, s) = G(x,\, \sigma s)
  \qquad \text{for all } x \ge 0,\ \sigma, s > 0. \tag{6.1}
\]
\statusT\quad\emph{the ``is equivalent to'' is carried by two declarations: (A8)
$\Rightarrow$ identities is \texttt{covariance\_laplace}, and identities $\Rightarrow$ (A8) is
\texttt{isScaleCovariant\_of\_repr\_map} (\texttt{Covariance.lean}, fidelity review R18).}
\end{lemma}

\begin{proof}
\leanok
The compatibility $\dil_\sigma(\mu * f) = (\dil_\sigma\mu) * (\dil_\sigma f)$ turns (A8) into
the statement that $\dil_\sigma\mu_{x,y}$ and $\mu_{\Sact_\sigma x, \Sact_\sigma y}$ represent
the same operator, and the uniqueness clause of \ref{lem:convolution-representation} makes
that an identity of measures; this is the first equivalence. For the second,
$\widehat{\dil_\sigma\mu}(s) = \hat\mu(\sigma s)$, so taking $-\log$ gives
$g_{\Sact_\sigma x, \Sact_\sigma y}(s) = g_{x,y}(\sigma s)$. Equation (6.1) is the case
$x = 0$ of that identity together with $\Sact_\sigma 0 = 0$, which holds because an increasing
bijection of $[0,\infty)$ fixes $0$; $G$ is $g_{0,\cdot}$ by \ref{lem:additivity}.
\end{proof}

% draft: Lemma 6.2
\begin{lemma}[rigidity of the action]\label{lem:action-rigidity}
\uses{def:cascade-family,lem:covariance-laplace,cor:strict-monotonicity,lem:additivity}
\lean{Hemigroup.CascadeCore.action_rigidity}\leanok
Assume (A1)--(A8) and (ND). Then:
\begin{enumerate}
\item For each $\sigma$, $\Sact_\sigma$ is uniquely determined by (6.1).
\item $\Sact_\sigma \Sact_\tau = \Sact_{\sigma\tau}$ for all $\sigma, \tau > 0$, and
$\Sact_1 = \Id$.
\item $\sigma \mapsto \Sact_\sigma x$ is continuous and strictly increasing for each $x > 0$.
\item The action has no fixed point in $(0,\infty)$: if $\Sact_\sigma x^* = x^*$ for all
$\sigma$, then $x^* = 0$.
\end{enumerate}
\statusT\quad\emph{(three of the four clauses are injectivity of $G(\cdot,s)$ --- that is,
\ref{cor:strict-monotonicity} --- and nothing else. (1) \emph{is} that injectivity
(\texttt{eq\_of\_G\_eq}); (2) is two applications of (6.1) on top of it (\texttt{S\_comp},
\texttt{S\_one}); (4) adds only that $G(x,\cdot)$ is continuous at the origin with value $0$,
which turns ``constant on $(0,\infty)$'' into ``identically zero'' and so contradicts
\ref{cor:strict-monotonicity} above $x = 0$ (\texttt{eq\_zero\_of\_fixed}).}

\emph{What (3) needed, and the development did not have, was strict monotonicity in the
\emph{other} variable: the proof below reads it off the sign of a Bernstein function's
derivative. The substitute is \texttt{laplace\_strictAnti}, one line of measure theory --- the
support of $e^{-st} - e^{-s't}$ is exactly $\{t \ne 0\}$, and that set carries positive mass
precisely because the kernel is not $\delta_0$, which is (ND). No derivative and no Bernstein
class appears. With it, (3)'s continuity half does not need the inverse function the proof below
invokes: \texttt{StrictMonoOn.continuousAt\_of\_exists\_between} asks only for orbit points
squeezed between a given bound and $\Sact_{\sigma_0} x$, and continuity of
$\sigma \mapsto G(x,\sigma)$ supplies them. That also avoids having to show the range of
$G(\cdot,s_0)$ is a neighbourhood, which is the step the phrase ``on its range'' hides.)}
\end{lemma}

\begin{proof}
\leanok
\emph{(1) Uniqueness.} If $\Sact$ and $\Sact'$ both satisfy (6.1) then
$G(\Sact_\sigma x, s) = G(x,\sigma s) = G(\Sact'_\sigma x, s)$ for every $s > 0$. By
\ref{cor:strict-monotonicity}, $G(\cdot, s)$ is strictly increasing, hence injective, so
$\Sact_\sigma x = \Sact'_\sigma x$.

\emph{(2) The group law.} Applying (6.1) twice,
\[
  G(\Sact_\sigma \Sact_\tau x, s) = G(\Sact_\tau x, \sigma s) = G(x, \sigma\tau s)
  = G(\Sact_{\sigma\tau} x, s),
\]
and injectivity of $G(\cdot,s)$ gives $\Sact_\sigma \Sact_\tau = \Sact_{\sigma\tau}$. Taking
$\sigma = 1$ in (6.1) gives $G(\Sact_1 x, s) = G(x,s)$, whence $\Sact_1 = \Id$.

\emph{(3) Continuity and strict monotonicity in $\sigma$.} Fix $x > 0$ and any $s_0 > 0$. By
(6.1), $G(\Sact_\sigma x, s_0) = G(x, \sigma s_0) = g_{0,x}(\sigma s_0)$.

The map $\sigma \mapsto g_{0,x}(\sigma s_0)$ is continuous, and it is \emph{strictly}
increasing: $g_{0,x}$ is a nonzero member of $\BFz$ by \ref{cor:strict-monotonicity}, and a
nonzero member of $\BFz$ is strictly increasing, because its derivative is completely monotone
and not identically zero, hence strictly positive on $(0,\infty)$ --- a completely monotone
function vanishing at one point vanishes beyond it, and vanishing on an interval forces the
representing measure, and so the function, to be zero.

The map $G(\cdot, s_0)$ is continuous (\ref{lem:additivity}) and strictly increasing
(\ref{cor:strict-monotonicity}), so it has a continuous strictly increasing inverse on its
range. Composing,
$\Sact_\sigma x = G(\cdot,s_0)^{-1}\bigl(g_{0,x}(\sigma s_0)\bigr)$ is continuous and strictly
increasing in $\sigma$.

\emph{(4) No interior fixed point.} Suppose $\Sact_\sigma x^* = x^*$ for all $\sigma > 0$.
Then (6.1) gives $G(x^*, s) = G(x^*, \sigma s)$ for all $\sigma, s > 0$, so $G(x^*, \cdot)$ is
constant on $(0,\infty)$; its value is $G(x^*, \zp) = 0$ by \ref{lem:additivity}. Hence
$g_{0,x^*} \equiv 0$, i.e.\ $\mu_{0,x^*} = \delta_0$ and $\Phi_{0,x^*} = \Id$. By (ND) this
forces $x^* = 0$.
\end{proof}

% draft: Proposition 6.3
\begin{proposition}[orbit coordinate; the gauge]\label{prop:canonical-gauge}
\uses{lem:action-rigidity,lem:covariance-laplace,cor:strict-monotonicity,thm:increments-bernstein,def:bernstein-function,def:levy-exponent}
\lean{Hemigroup.CascadeCore.canonical_gauge}\leanok
Assume (A1)--(A8) and (ND). Then $c(\sigma) := \Sact_\sigma 1$ is a continuous strictly
increasing bijection of $(0,\infty)$ onto $(0,\infty)$ with $c(1) = 1$. Consequently
\[
  \chi := c^{-1} : (0,\infty) \to (0,\infty), \qquad \chi(0) := 0,
\]
is an increasing bijection --- \emph{the canonical gauge} --- satisfying
$\chi(\Sact_\sigma x) = \sigma\,\chi(x)$, and
\[
  G(x, s) = F\bigl(\chi(x)\, s\bigr), \qquad
  F := G(1, \cdot) \in \BFz,\ F \not\equiv 0. \tag{6.2}
\]
Equivalently $F \in \LE$, with a drift $b_0 \ge 0$ and a Lévy measure $\nu$ on $(0,\infty)$.
\statusT\quad\emph{(the $\LE$ reading is the same statement, by \ref{prop:bernstein-toolbox}(3),
and is the one \ref{thm:main-analysis} needs: it is what feeds
\ref{lem:selfdecomposable-exponents} and, through it, the form (7.1). $F$'s membership comes
straight from \ref{thm:increments-bernstein}, so nothing new is asserted here --- the class is
carried along, not re-established. \textbf{Formalised}, and the interface question the
formalisation note flagged (M5 of the ladder) is settled the way it warned it had to be:
\texttt{canonical\_gauge} asserts $\chi$ together with the conjugation
$\chi(\Sact_\sigma x) = \sigma\,\chi(x)$ and the similarity form $G(x,s) = G(1,\chi(x)\,s)$,
not $\chi$ alone. That bundle is what \ref{prop:main-uniqueness} declined to guess at.}

\emph{Only surjectivity of $c$ had to be proved; continuity and strict monotonicity are
\ref{lem:action-rigidity}(3) and $c(1) = 1$ is the group law, so the injectivity argument below
--- the iteration along $\rho^{\pm n}$ --- is not needed in Lean at all. Surjectivity is done
without limits. The orbit $R = c((0,\infty))$ is invariant under every $\Sact_\rho$, which
\emph{is} the group law; were it bounded above, its supremum would be fixed by the action ---
nothing on the orbit can be carried above it, and nothing below it either, since
$\Sact_{\rho^{-1}}$ would then carry it down --- and \ref{lem:action-rigidity}(4) forbids a fixed
point above the origin, while $1 \in R$. Bounded below is the same argument read downwards, where
a positive infimum is itself the contradiction. The intermediate value theorem then finishes.)}
\end{proposition}

\begin{proof}
\leanok
\emph{$c$ is continuous} by \ref{lem:action-rigidity}(3), and $c(1) = \Sact_1 1 = 1$ by
\ref{lem:action-rigidity}(2).

\emph{$c$ is injective.} Suppose $c(\sigma) = c(\tau)$ with $\rho := \sigma/\tau \ne 1$. The
group law \ref{lem:action-rigidity}(2) gives $\Sact_\rho 1 = 1$, so by (6.1)
\[
  F(s) = G(1,s) = G(\Sact_\rho 1, s) = G(1, \rho s) = F(\rho s)
  \qquad \text{for all } s > 0 .
\]
Iterating along $\rho^{\pm n}$ gives $F(s) = F(\rho^n s)$ for every $n \in \ZZ$. Replacing
$\rho$ by $\rho^{-1}$ if necessary we may take $\rho < 1$, and then $\rho^n s \to 0$, so
$F(s) = F(\zp) = 0$ for every $s$ by \ref{lem:additivity}. Thus $F \equiv 0$, which
contradicts \ref{cor:strict-monotonicity} under (ND).

\emph{$c$ is increasing.} Being continuous and injective on the connected set $(0,\infty)$,
$c$ is strictly monotone, so it suffices to exclude the decreasing case. Suppose $c$ were
decreasing and fix $\sigma > 1$, so that $\Sact_\sigma 1 < 1$. Since $G(\cdot, s)$ is strictly
increasing (\ref{cor:strict-monotonicity}),
\[
  G(\Sact_\sigma 1, s) < G(1,s) = F(s) \qquad \text{for every } s > 0 .
\]
But (6.1) evaluates the left-hand side as $G(1,\sigma s) = F(\sigma s)$, and $F$ is
nondecreasing, being a Bernstein function, so $F(\sigma s) \ge F(s)$. The two are
incompatible, so $c$ is increasing.

\emph{$c$ is onto $(0,\infty)$.} Its limits as $\sigma \to \zp$ and $\sigma \to \infty$ exist
in $[0,\infty]$ by monotonicity, and each is a fixed point of the action or a boundary point;
by \ref{lem:action-rigidity}(4) there is no fixed point in $(0,\infty)$, so the limits are $0$
and $\infty$. With continuity, $c$ is a bijection onto $(0,\infty)$.

\emph{The gauge relation.} By the group law, $\Sact_\sigma x = \Sact_\sigma \Sact_{\chi(x)} 1
= \Sact_{\sigma\chi(x)} 1$, so $\chi(\Sact_\sigma x) = \sigma\chi(x)$.

\emph{The similarity form.} For $x > 0$, writing $x = \Sact_{\chi(x)} 1$ and using (6.1),
\[
  G(x, s) = G\bigl(\Sact_{\chi(x)} 1,\, s\bigr) = G\bigl(1,\, \chi(x)\, s\bigr)
  = F\bigl(\chi(x)\,s\bigr),
\]
and $F = g_{0,1} \in \BFz$ by \ref{thm:increments-bernstein}, with $F \not\equiv 0$ by (ND)
through \ref{cor:strict-monotonicity}. The case $x = 0$ is $G(0,\cdot) = 0 = F(0)$.
\end{proof}

% draft: Remark 6.4
\begin{remark}[gauge freedom]\label{rem:gauge-freedom}
\uses{prop:canonical-gauge}
In the hemigroup setting the scale coordinate carries no intrinsic parametrization: replacing
$x$ by any increasing bijection of $[0,\infty)$ yields an equivalent family.
\ref{prop:canonical-gauge} fixes the \emph{canonical gauge} $\tilde x := \chi(x)$, in which
$\Sact_\sigma$ acts by multiplication, scale carries the dimension of time, and
$G(x,s) = F(\tilde x\, s)$.

The powers $\tilde x \mapsto \tilde x^{\,\alpha}$ parametrize the other homogeneous gauges;
the \emph{parabolic gauge} $\alpha = 1/2$ is the natural one for the diffusive memory-line
realization and is used in the sequel.
\end{remark}
